\chapter{Технологический раздел}

\section{Выбор языка и среды программирования}

Модуль реализован на языке C с использованием штатного диалекта ядра GNU11. Собственная вычислительная логика использует конструкции C99, однако заголовки ядра требуют расширений GNU. Поэтому режим строгого ISO C99 для всего модуля не задаётся~\cite{kernellang,gcc}.

Средства разработки включают редактор, GCC, GNU Make и подготовленные заголовки целевого ядра. Kbuild согласует конфигурацию, архитектурные заголовки и экспортируемые символы~\cite{kbuild}. Административное средство написано на Python 3; стандартная библиотека обеспечивает работу с sysfs, JSON, файловыми блокировками и системными командами~\cite{python}. Python не участвует в обработке движения. Проект целиком размещается в целевой Linux-системе, где выполняются сборка и эксплуатация.

\section{Структуры и точки входа}

Контекст устройства приведён в листинге~\ref{lst:device-context}. Поле \path{udev} удерживает ссылку на USB-устройство, \path{interface} указывает интерфейс. Запрос \path{urb} переиспользуется; \path{buffer} и \path{dma} описывают буфер для процессора и DMA. Состояние пера и разрешение отправки защищены спин-блокировкой.

\par\noindent\begin{minipage}{\linewidth}
\lstinputlisting[language=C,basicstyle=\footnotesize\ttfamily,caption={Контекст устройства},label={lst:device-context},firstline=68,lastline=81]{../src/main.c}
\end{minipage}\par

Флаг \path{running} означает разрешение отправки, а не постоянное наличие активного URB: во время восстановления запрос временно не отправлен. Точки входа связаны структурой из листинга~\ref{lst:entry-points}. Инициализация проверяет делитель и вызывает \path{usb_register()}, завершение --- \path{usb_deregister()}. Аппаратную совместимость ограничивают таблица идентификаторов и разборщик отчёта.

\par\noindent\begin{minipage}{\linewidth}
\lstinputlisting[language=C,basicstyle=\footnotesize\ttfamily,caption={Точки входа драйвера},label={lst:entry-points},firstline=322,lastline=332]{../src/main.c}
\end{minipage}\par

\section{Подготовка приёма и устройства ввода}

Функция \path{xppen_probe()} выбирает interrupt-конечную точку через \path{usb_find_int_in_endpoint()}. Размер буфера и интервал берутся из её дескриптора. Память контекста выделяет \path{kzalloc()}, запрос создаёт \path{usb_alloc_urb()}. Для буфера используется \path{usb_alloc_coherent()}. Эти операции выполняются с флагом \path{GFP_KERNEL}. Для готового DMA-отображения устанавливается \path{URB_NO_TRANSFER_DMA_MAP}~\cite{kernelusb}.

Устройство выделяется через \path{input_allocate_device()}. Его возможности ограничиваются \path{REL_X}, \path{REL_Y}, \path{BTN_LEFT} и \path{BTN_RIGHT}; затем вызывается \path{input_register_device()}~\cite{kernelinput}. Функция \path{xppen_start()} устанавливает Report Protocol стандартным запросом HID, не изменяя дескриптор отчёта~\cite{usb_hid}. При успехе сбрасывается начало отсчёта и запускается приём. При ошибке освобождаются уже полученные ресурсы с учётом регистрации устройства ввода.

\section{Преобразование и остановка обмена}

Функция \path{pen_decode()} получает состояние пера, буфер, фактическую длину, делитель и выходной кадр. Проверки предшествуют изменению сохранённого отсчёта; конкретная раскладка пакета описана в исследовательском разделе. Расчёт приращений приведён в листинге~\ref{lst:delta}. Приведение делителя к знаковому типу сохраняет знаковое деление отрицательных смещений.

\par\noindent\begin{minipage}{\linewidth}
\lstinputlisting[language=C,basicstyle=\footnotesize\ttfamily,caption={Расчёт приращений и остатков},label={lst:delta},firstline=50,lastline=61]{../src/main.c}
\end{minipage}\par

Делитель не меняется во время работы, поэтому остатки сохраняют один масштаб. Обработчик \path{xppen_complete()} использует \path{urb->actual_length}, передаёт допустимый кадр через функции \path{input_report_rel()} и \path{input_report_key()}, завершая его \path{input_sync()}. Отправка выполняется в \path{xppen_submit_locked()} под блокировкой с \path{GFP_ATOMIC}, без выделения памяти для очередного пакета.

Общая процедура остановки приведена в листинге~\ref{lst:stop}. Разрешение отправки снимается до ожидания URB и рабочей функции. После ожидания освобождаются кнопки; при отключении дополнительно снимается регистрация устройства ввода и возвращаются ресурсы.

\par\noindent\begin{minipage}{\linewidth}
\lstinputlisting[language=C,basicstyle=\footnotesize\ttfamily,caption={Остановка приёма},label={lst:stop},firstline=190,lastline=200]{../src/main.c}
\end{minipage}\par

Статусы отмены прекращают цикл. При \path{-EPIPE} рабочая функция вызывает \path{usb_clear_halt()} вне спин-блокировки и повторно проверяет разрешение приёма. Ошибка восстановления освобождает кнопки. Для остальных ошибок выполняется повторная отправка с ограниченной по частоте диагностикой.

\section{Управление и сборка}

Режим работы Makefile определяется переменной \path{KERNELRELEASE}. Внешний вызов передаёт каталог проекта параметром \path{M}; внутри Kbuild объявляется состав модуля~\cite{kbuild}.

\par\noindent\begin{minipage}{\linewidth}
\lstinputlisting[language=make,basicstyle=\footnotesize\ttfamily,caption={Настройка сборки модуля},label={lst:make-build},firstline=1,lastline=13]{../Makefile}
\end{minipage}\par

Переменная \path{KDIR} выбирает каталог сборки ядра, \path{DIVISOR} --- чувствительность, \path{PYTHON} --- интерпретатор. Предупреждения \path{W=1} включены в команду сборки. Загрузка, привязка, выгрузка и очистка реализованы целями из листинга~\ref{lst:make-control}.

\par\noindent\begin{minipage}{\linewidth}
\lstinputlisting[language=make,basicstyle=\footnotesize\ttfamily,caption={Команды управления модулем},label={lst:make-control},firstline=16,lastline=27]{../Makefile}
\end{minipage}\par

Средство \path{scripts/manage.py} проверяет Linux и права администратора, защищает каталог состояния и сериализует команды файловой блокировкой. Привязки сохраняются до изменения sysfs; повторная привязка не перезаписывает имя штатного драйвера. При восстановлении заново проверяются идентификаторы устройств, а неудачные операции сохраняются для повторения. Автоматическая перепривязка правилами udev не устанавливается: переключение выполняется явной командой администратора~\cite{udev}. Очистка артефактов не требует заголовков ядра.

Таким образом, реализованы асинхронный USB-приём, преобразование событий и согласованная остановка. Административные операции вынесены из пути обработки ввода и обеспечивают возврат к штатной обработке.

\clearpage
